\documentclass[11pt]{letter}
\usepackage[margin=1in]{geometry}

\signature{Decory Edwards}

\begin{document}

\begin{letter}{Search Committee \\
MIT Center for Energy and Environmental Policy Research (CEEPR) \\
MIT Sloan School of Management}

\opening{Dear Professor Knittel and Members of the Search Committee,}

I am writing to apply for the Postdoctoral Associate position at the MIT Center for Energy and Environmental Policy Research (CEEPR). I am currently a Ph.D.\ candidate in Economics at Johns Hopkins University and expect to complete my degree in May 2026. My research focuses on quantitative economic modeling and empirical analysis of markets and policy, and I would be excited to contribute to CEEPR’s work on the economics of the energy transition, particularly the supply chains and market dynamics surrounding critical minerals.

My research examines how economic incentives and market frictions shape long-run economic outcomes using a combination of structural modeling and empirical analysis. In my job market paper, I develop a heterogeneous-agent life-cycle model to study the distributional dynamics of wealth accumulation, incorporating differences in returns to wealth and income risk. The project combines economic theory with empirical calibration using microdata from the Survey of Consumer Finances and requires extensive data construction, econometric analysis, and computational modeling using Python and Stata. Through this work I have developed strong experience building and estimating quantitative models and working with large datasets to answer policy-relevant economic questions.

I am particularly interested in the economic forces shaping the energy transition and the markets for the inputs required to support it. The rapid expansion of electric vehicles, renewable energy infrastructure, and energy storage technologies has increased demand for minerals such as lithium, nickel, and cobalt, while supply remains constrained by geological, technological, and geopolitical factors. I am excited by the opportunity to help develop empirical models that incorporate both economic incentives and the physical constraints of mining operations in order to analyze how policies such as EV subsidies, trade restrictions, or domestic content requirements influence global mineral markets.

My methodological training aligns closely with the skills required for this project. I have extensive experience with applied econometrics and quantitative modeling, including panel data analysis, structural modeling, and computational economic simulations. I regularly work with statistical software such as Python and Stata, and I am comfortable building empirical models that combine economic theory with real-world data. I also value collaborative research environments and would welcome the opportunity to work with Professor Knittel and other researchers at CEEPR on interdisciplinary questions at the intersection of energy markets, trade, and public policy.

MIT CEEPR’s work on energy and environmental policy has had an important impact on both academic research and policy discussions, and I would be excited to contribute to this effort. I am available to begin the position in Summer or Fall 2026 and would welcome the opportunity to discuss how my research background and quantitative skills could contribute to this project.

Thank you for your time and consideration.

\closing{Sincerely,}

\end{letter}

\end{document}